% Originally: nothing. Chapter 24 had no exercises at all and therefore no solutions section;
% both are new on 2026-08-09. The chapter is the most abstract in the document and had the least
% for a reader to do, which is the wrong way round.
% Generator: Claude Opus 5 (High)

\section{Solutions}
\label{sec:differential_forms_solutions}

\begin{exercisesolution}[Solution to \cref{ex:ai_three_shadows}]
\begin{enumerate}[label=\textbf{(\alph*)}]
  \item \textbf{The three shadows.} Property \textbf{(d)} says each is the determinant of the
    $2\times2$ block of $(v \mid w)$ cut out by the two indices, so with $v = (1,2,0)$ and
    $w = (0,1,3)$:
    \begin{align*}
      dx^1\wedge dx^2(v,w) &= \det\begin{pmatrix}1 & 0\\ 2 & 1\end{pmatrix} = 1, \\
      dx^1\wedge dx^3(v,w) &= \det\begin{pmatrix}1 & 0\\ 0 & 3\end{pmatrix} = 3, \\
      dx^2\wedge dx^3(v,w) &= \det\begin{pmatrix}2 & 1\\ 0 & 3\end{pmatrix} = 6.
    \end{align*}
  \item \textbf{The cross product.}
    \[v \times w = \big(2\cdot3 - 0\cdot1,\ 0\cdot0 - 1\cdot3,\ 1\cdot1 - 2\cdot0\big)
      = (6,\,-3,\,1).\]
  \item \textbf{The correspondence.} Component by component,
    \begin{align*}
      (v\times w)_1 &= 6 = dx^2\wedge dx^3(v,w), \\
      (v\times w)_2 &= -3 = -\,dx^1\wedge dx^3(v,w), \\
      (v\times w)_3 &= 1 = dx^1\wedge dx^2(v,w).
    \end{align*}
    So the $1$st and $3$rd match on the nose and the $2$nd needs a sign flip. The pattern becomes
    uniform once that middle shadow is written in the cyclic order $3,1$ instead of the increasing
    order $1,3$: since $dx^3\wedge dx^1 = -\,dx^1\wedge dx^3$, one has
    \[v\times w = \big(dx^2\wedge dx^3(v,w),\ dx^3\wedge dx^1(v,w),\ dx^1\wedge dx^2(v,w)\big).\]
\end{enumerate}

\begin{remark}[The cross product is a $2$-form in disguise]
\label{rem:cross_product_is_a_two_form}
This is not a coincidence about these two vectors. Comparing with example item \textbf{(d)} of
\cref{sec:differential_forms}, the identity says exactly that
$\omega^2_{v\times w} = dx^1\wedge dx^2$-and-friends evaluated on $(v,w)$, i.e.\ that the cross
product \emph{is} the $2$-form $(v,w) \mapsto \det(\,\cdot\mid v\mid w)$ read as a vector through
the $\mathbb{R}^3$ inner product.

Two things fall out. The alternating sign that everyone memorises in the cofactor expansion of
$v\times w$ is the same sign as $dx^3\wedge dx^1 = -\,dx^1\wedge dx^3$, so it is antisymmetry
rather than a mnemonic. And the cross product exists only in $\mathbb{R}^3$ because that is the
only dimension where $\binom{n}{2} = n$, i.e.\ where there are exactly as many $2$-form shadows
as coordinates to store them in. In $\mathbb{R}^4$ there are six shadows and four slots, and no
cross product.
\end{remark}
\end{exercisesolution}

\begin{exercisesolution}[Solution to \cref{ex:ai_wedge_square}]
\begin{enumerate}[label=\textbf{(\alph*)}]
  \item Expand by bilinearity, writing $dx^{ij} := dx^i\wedge dx^j$:
    \[\alpha\wedge\alpha = dx^{12}\wedge dx^{12} + dx^{12}\wedge dx^{34} + dx^{34}\wedge dx^{12}
      + dx^{34}\wedge dx^{34}.\]
    The first and last vanish, since each repeats a factor and $dx^i\wedge dx^i = 0$. For the two
    middle terms, moving the block $dx^{34}$ past the block $dx^{12}$ means moving two $1$-forms
    past two $1$-forms, which costs $(-1)^{2\cdot2} = +1$, so $dx^{34}\wedge dx^{12} =
    dx^{12}\wedge dx^{34}$ and the two middle terms add rather than cancel:
    \[\alpha\wedge\alpha = 2\,dx^1\wedge dx^2\wedge dx^3\wedge dx^4 \neq 0.\]
  \item The argument for a $1$-form is $\beta\wedge\beta = -\beta\wedge\beta$, which forces
    $\beta\wedge\beta = 0$. That commutation sign comes from the graded rule: for a $k$-form
    $\alpha$ and an $m$-form $\theta$,
    \[\alpha\wedge\theta = (-1)^{km}\,\theta\wedge\alpha.\]
    With $\theta = \alpha$ and $m = k$ this reads
    $\alpha\wedge\alpha = (-1)^{k^2}\alpha\wedge\alpha$, and $(-1)^{k^2} = (-1)^k$. So the
    argument gives $\alpha\wedge\alpha = 0$ exactly when $k$ is \textbf{odd}, and says nothing at
    all when $k$ is even. The $2$-form above is the standard witness that \qt{nothing at all} is
    the correct reading.
\end{enumerate}
\end{exercisesolution}

\begin{exercisesolution}[Solution to \cref{ex:ai_angle_form_closed}]
\begin{enumerate}[label=\textbf{(\alph*)}]
  \item Write $\omega = P\,dx + Q\,dy$ with $P = \dfrac{-y}{x^2+y^2}$ and
    $Q = \dfrac{x}{x^2+y^2}$. Then $d\omega = (\partial_xQ - \partial_yP)\,dx\wedge dy$, and by
    the quotient rule
    \[\partial_xQ = \frac{(x^2+y^2) - x\cdot 2x}{(x^2+y^2)^2} = \frac{y^2-x^2}{(x^2+y^2)^2},
      \qquad
      \partial_yP = \frac{-(x^2+y^2) + y\cdot 2y}{(x^2+y^2)^2} = \frac{y^2-x^2}{(x^2+y^2)^2}.\]
    They are equal, so $d\omega = 0$ on all of $U$.
  \item Along $\gamma(t) = (\cos t,\sin t)$ the denominator is $1$, and
    $-y\,dx + x\,dy$ pulls back to $(-\sin t)(-\sin t)\,dt + (\cos t)(\cos t)\,dt = dt$. Hence
    \[\int_\gamma \omega = \int_0^{2\pi} dt = 2\pi \neq 0.\]
    Under the identification $F \leftrightarrow \omega_F$, an exact form $\omega = df$
    corresponds to a conservative field with potential $f$, and
    \cref{prop:work_difference_potential} makes its integral around any closed loop
    $f(\gamma(1)) - f(\gamma(0)) = 0$. Since $2\pi \neq 0$, no such $f$ exists: $\omega$ is
    \textbf{closed but not exact}.
\end{enumerate}

\begin{remark}[What the counterexample is really about]
\label{rem:angle_form_hole}
The obstruction is not analytic but topological. On any simply connected subset of $U$ --- the
plane slit along a ray, say --- the form \emph{is} exact, with $f$ the polar angle $\theta$,
which is why $\omega$ is called the \newterm{angle form}. What fails on all of
$\mathbb{R}^2\setminus\{0\}$ is that $\theta$ cannot be defined continuously all the way round:
it gains $2\pi$ per circuit, and $2\pi$ is exactly the integral computed above.

So \qt{closed} and \qt{exact} differ precisely by the holes in the domain. That is the content of
the Poincaré lemma (\cref{thm:poincare_lemma}), which says the two coincide on star-shaped
domains, and of \cref{cor:potentials_simply_connected}. This one-line example is the reason those
theorems need a hypothesis on the domain at all.
\end{remark}
\end{exercisesolution}

\begin{exercisesolution}[Solution to \cref{ex:ai_exterior_derivative_drill}]
\begin{enumerate}[label=\textbf{(\alph*)}]
  \item Using $d\omega = \sum_I d\omega_I\wedge dx^I$ and discarding every wedge with a
    repeated factor,
    \begin{align*}
    d\omega &= d(xy)\wedge dx + d(yz)\wedge dy + d(zx)\wedge dz \\
      &= (y\,dx + x\,dy)\wedge dx + (z\,dy + y\,dz)\wedge dy + (z\,dx + x\,dz)\wedge dz \\
      &= x\,dy\wedge dx + y\,dz\wedge dy + z\,dx\wedge dz \\
      &= -x\,dx\wedge dy - y\,dy\wedge dz - z\,dz\wedge dx.
    \end{align*}
    By item \textbf{(b)} of \cref{ex:exterior_derivative_gradient_curl} this is
    $\omega^2_{\curl F}$ for $F = (xy,\,yz,\,zx)$, so the computation is
    $\curl F = (-y,\,-z,\,-x)$. (Check directly:
    $\curl F = (\partial_y(zx) - \partial_z(yz),\ \partial_z(xy) - \partial_x(zx),\
    \partial_x(yz) - \partial_y(xy)) = (-y,\,-z,\,-x)$.)
  \item Only the $dx_1$ part of $d(x_1x_2x_3)$ survives, the other two repeating a factor
    already present:
    \[d\omega = d(x_1x_2x_3)\wedge dx_2\wedge dx_3 = x_2x_3\,dx_1\wedge dx_2\wedge dx_3.\]
  \item By property \textbf{(b)} of \cref{thm:exterior_derivative} with $k = 0$, and then
    $d(dg) = 0$ by property \textbf{(c)},
    \[d(f\,dg) = df\wedge dg + f\,d(dg) = df\wedge dg.\]
    Swapping the roles of $f$ and $g$ gives $d(g\,df) = dg\wedge df = -\,df\wedge dg$, since both
    are $1$-forms. Adding,
    \[d(f\,dg + g\,df) = df\wedge dg - df\wedge dg = 0.\]
    This is the Leibniz rule in disguise: $f\,dg + g\,df = d(fg)$, and $d\circ d = 0$.
\end{enumerate}
\end{exercisesolution}
